\apartado{2.6}{Teorema de Bayes}{climax}
\bajada{Dar vuelta una probabilidad condicional. Es el apartado que explica por
qué un test que casi nunca falla puede equivocarse en la mitad de sus alarmas —
y por qué el mismo instrumento sirve en un servicio y no sirve en otro.}

\definicion{Valor predictivo positivo}{%
  La probabilidad de tener la condición \emph{dado que} el test dio positivo.
  Es la pregunta del paciente, y no es la misma que la sensibilidad, que es la
  pregunta del instrumento.}

\minihistoria{Los médicos también se equivocan en esto}{%
  En los años setenta, Kahneman y Tversky le plantearon el problema a médicos,
  psicólogos y estadísticos entrenados: una enfermedad afecta al 1\,\% de la
  población, un test la detecta con 95\,\% de acierto, una persona da positivo.
  ¿Qué probabilidad tiene de estar enferma?\par
  La respuesta mayoritaria fue 95\,\%. La correcta rondaba el 16\,\%. El error
  no era de aritmética: al ver el resultado del test, las personas se olvidaban
  de un dato que tenían delante — que la enfermedad es rara.}

\ejercicio{0}{Define con tus palabras qué es la \emph{prevalencia} de un
  trastorno. ¿De qué depende: del instrumento que se use, o de la población que
  se estudie?}{}
\renglones[3]
\respuesta{La proporción de la población que tiene el trastorno. Depende de la
población, no del instrumento.}
\solucion{%
  Es la \textbf{proporción de la población que tiene la condición} en un
  momento dado.
  \par
  Depende de la \textbf{población}, no del instrumento. La misma prueba
  aplicada en una guardia psiquiátrica y en una feria de la salud da resultados
  de utilidad muy distinta, porque la prevalencia es distinta aunque el
  instrumento sea idéntico.
  \par
  Ésa es toda la idea del apartado, y la razón por la que no existe «el valor
  predictivo de un test» sin decir en qué población.}

\ejercicio{1}{Un tamizaje de riesgo suicida se aplica en una población donde la
  condición afecta al 1\,\% de las personas. El test tiene 90\,\% de
  sensibilidad y 90\,\% de especificidad.
  \textbf{Antes de calcular, escribe qué probabilidad crees que tiene de estar
  en riesgo alguien que dio positivo:} \underline{\hspace{2cm}}
  \par\vspace{0.2em}
  Lo vas a comprobar en el ejercicio 4.}{Predice primero}
\renglones[3]
\respuesta{Es 8,3\,\%. Casi todo el mundo estima cerca del 90\,\%.}
\solucion{%
  La respuesta es \textbf{8,3\,\%}, y se calcula en el ejercicio 4.
  \par
  La estimación intuitiva casi siempre ronda el 90\,\%, porque uno toma el
  acierto del test como si fuera la respuesta. Es el mismo error que cometieron
  los médicos del estudio de Kahneman y Tversky.
  \par
  La distancia entre tu estimación y el 8,3\,\% es lo que este apartado
  corrige.}
\enclase{Anota las estimaciones antes de seguir. Si el curso dice 90\,\%, está
en buena compañía: lo mismo respondieron profesionales con años de práctica.}

\ejercicio{1}{De cada 1000 personas de cierta población, 10 tienen el
  trastorno.
  \begin{enumerate}[label=(\alph*),topsep=2pt,itemsep=1pt]
    \item ¿Cuántas no lo tienen?
    \item ¿Cuál es la prevalencia, en porcentaje?
  \end{enumerate}}{}
\renglones[3]
\respuesta{(a) 990. (b) 1\,\%.}
\solucion{%
  (a) $1000-10=990$. \quad (b) $\dfrac{10}{1000}=0{,}01$, o sea 1\,\%.
  \par
  Trabajar con \emph{frecuencias naturales} («10 de cada 1000») en vez de
  porcentajes reduce muchísimo los errores. Es un resultado experimental
  conocido, no una preferencia de estilo: la misma persona que falla el
  problema en porcentajes lo resuelve bien en frecuencias.}

\ejercicio{2}{En cada frase, identifica si se habla de $P(+\mid E)$ —dar
  positivo dado que se tiene el trastorno— o de $P(E\mid +)$ —tener el
  trastorno dado que se dio positivo—.
  \begin{enumerate}[label=(\alph*),topsep=2pt,itemsep=1pt]
    \item «De los pacientes con depresión, el test detecta al 88\,\%.»
    \item «De los que dieron positivo, la mitad tenía el trastorno.»
    \item «El instrumento tiene una sensibilidad del 88\,\%.»
  \end{enumerate}
  Después escribe cuál es el truco de lectura que te sirvió.}{Dos casos casi iguales}
\renglones[5]
\respuesta{(a) $P(+\mid E)$; (b) $P(E\mid +)$; (c) $P(+\mid E)$. El truco: lo
que viene después de «de los\ldots» es el denominador.}
\solucion{%
  (a) $P(+\mid E)$ — parte de los que tienen el trastorno.\par
  (b) $P(E\mid +)$ — parte de los que dieron positivo.\par
  (c) $P(+\mid E)$ — «sensibilidad» es otro nombre para lo mismo.
  \par
  \textbf{El truco de lectura}: busca la expresión «de los\ldots». Lo que viene
  inmediatamente después es el \emph{denominador}, o sea el grupo del que ya
  sabemos algo. Ahí está la condición.
  \par
  Las frases (a) y (c) dicen exactamente lo mismo con palabras distintas, y la
  (b) dice algo completamente diferente. Ésa es toda la dificultad del
  apartado.}

\ejercicio{2}{Retoma el tamizaje del ejercicio 2: prevalencia 1\,\%,
  sensibilidad 90\,\%, especificidad 90\,\%. Usando \textbf{frecuencias
  naturales} sobre 1000 personas, completa:
  \begin{enumerate}[label=(\alph*),topsep=2pt,itemsep=1pt]
    \item ¿Cuántas tienen el trastorno? ¿Cuántas de ellas dan positivo?
    \item ¿Cuántas no lo tienen? ¿Cuántas de ellas dan positivo igual?
    \item ¿Cuántos positivos hay en total? ¿Qué proporción de ellos lo tiene?
  \end{enumerate}}{}
\renglones[7]
\respuesta{(a) 10 con el trastorno, 9 positivas. (b) 990 sin él, 99 positivas.
(c) 108 positivos, de los cuales 9 lo tienen: 8,3\,\%.}
\solucion{%
  (a) Con el trastorno: 1\,\% de 1000 es 10. De ellas dan positivo el 90\,\%:
  \textbf{9}.
  \par
  (b) Sin el trastorno: 990. Con 90\,\% de especificidad, el 10\,\% da positivo
  igual: 99 falsas alarmas.
  \par
  (c) Positivos totales: $9+99=108$. De ellos lo tienen 9:
  \[ \frac{9}{108}\approx 0{,}083 \]
  o sea 8,3\,\%.
  \par
  \textbf{Un test con 90\,\% de acierto, y sin embargo 9 de cada 10 alarmas son
  falsas.} No es un defecto del test: es que hay 99 veces más gente sana que
  enferma, así que aun equivocándose poco con los sanos, en números absolutos
  se equivoca muchísimo.
  \par
  Compara con tu estimación del ejercicio 2.}
\enclase{Dibuja el árbol en el pizarrón con las 1000 personas. Ver los 99
falsos positivos al lado de los 9 verdaderos convence más que la fórmula.}

\ejercicio{3}{Rehaz el cálculo anterior con la \textbf{fórmula de Bayes} y
  comprueba que da lo mismo:
  \[ P(E\mid +) = \frac{P(+\mid E)\cdot P(E)}
       {P(+\mid E)\cdot P(E) + P(+\mid \bar{E})\cdot P(\bar{E})} \]}{}
\renglones[5]
\respuesta{$(0{,}9\times 0{,}01)/(0{,}9\times 0{,}01+0{,}1\times 0{,}99)
\approx 0{,}083$.}
\solucion{%
  \[ P(E\mid +) = \frac{0{,}90\times 0{,}01}
     {0{,}90\times 0{,}01 + 0{,}10\times 0{,}99}
     = \frac{0{,}009}{0{,}009+0{,}099}=\frac{0{,}009}{0{,}108}\approx 0{,}083 \]
  \par
  Mismo resultado. Y fíjate que $0{,}009$ y $0{,}099$ son los 9 y los 99 del
  ejercicio anterior divididos por 1000: \textbf{la fórmula y las frecuencias
  naturales son literalmente la misma cuenta}, escrita con proporciones en vez
  de personas.
  \par
  Conviene saber las dos: las frecuencias naturales para entender y para
  explicarle a un paciente, la fórmula para cuando no tienes la tabla pero sí
  los tres números.}

\ejercicio{3}{Aplica Bayes al servicio real: sensibilidad $0{,}88$,
  especificidad $0{,}88$, prevalencia $0{,}125$.
  \begin{enumerate}[label=(\alph*),topsep=2pt,itemsep=1pt]
    \item Calcula la probabilidad de tener el trastorno dado que se dio
      positivo.
    \item Compárala con el $22/43$ que sale de contar la tabla del apartado 2.3.
    \item Repite el cálculo suponiendo que el mismo instrumento se aplica en
      una guardia donde la prevalencia es del 50\,\%.
  \end{enumerate}}{}
\renglones[8]
\respuesta{(a) $\approx 0{,}512$. (b) Idéntico a $22/43$. (c) $0{,}88$.}
\solucion{%
  (a) \[ \frac{0{,}88\times 0{,}125}{0{,}88\times 0{,}125+0{,}12\times 0{,}875}
     = \frac{0{,}110}{0{,}215}\approx 0{,}512 \]
  \par
  (b) Contando la tabla: $\dfrac{22}{43}\approx 0{,}512$. \textbf{Idéntico}, y
  no es casualidad: $0{,}110\times 200=22$ y $0{,}215\times 200=43$.
  \par
  (c) Con prevalencia $0{,}50$:
  \[ \frac{0{,}88\times 0{,}50}{0{,}88\times 0{,}50+0{,}12\times 0{,}50}
     = \frac{0{,}44}{0{,}50}=0{,}88 \]
  \par
  \textbf{El valor predictivo pasó de 51,2\,\% a 88\,\% sin tocar el
  instrumento.} Cambió la población.
  \par
  Consecuencia práctica: cuando leas que un instrumento tiene «buen valor
  predictivo», pregunta en qué población se validó. Uno validado en un servicio
  de urgencias puede ser casi inútil como tamizaje poblacional, y el número
  publicado no te avisa.}

\ejercicio{4}{Una estudiante da positivo en el tamizaje del servicio y
  pregunta, angustiada, si eso significa que tiene depresión.
  \begin{enumerate}[label=(\alph*),topsep=2pt,itemsep=1pt]
    \item ¿Qué número le corresponde como respuesta, y por qué ése y no el
      88\,\%?
    \item Redacta lo que le dirías, sin usar la palabra «probabilidad».
    \item Si el servicio bajara el corte de 10 a 8, ¿el valor predictivo
      subiría o bajaría? Justifica sin calcular.
  \end{enumerate}}{}
\renglones[8]
\respuesta{(a) 51,2\,\%, el valor predictivo positivo. (b) Respuesta abierta.
(c) Bajaría.}
\solucion{%
  (a) Le corresponde el \textbf{valor predictivo positivo: 51,2\,\%}. El 88\,\%
  responde otra pregunta —«si tuviera depresión, ¿la detectaríamos?»— y ella ya
  sabe que dio positivo: su denominador son los 43 marcados, no los 25 con el
  trastorno.
  \par
  (b) Algo del tipo: «Este cuestionario está hecho para marcar de más, a
  propósito: preferimos conversar de más antes que dejar pasar a alguien. De
  cada dos personas que marcamos, aproximadamente una no tiene el trastorno. Lo
  que sigue es una entrevista para mirarlo con calma, no un diagnóstico.»
  \par
  (c) \textbf{Bajaría.} Al bajar el corte se marca a más personas; se captan
  algunos casos que antes se escapaban, pero la mayoría de los nuevos marcados
  son sanos, simplemente porque los sanos son muchos más. El numerador crece
  poco y el denominador crece mucho. Se comprueba numéricamente en 2.9.
  \par
  \textbf{Adónde lleva esto}: el mismo error de invertir una condicional
  reaparece en investigación con el valor $p$. Un $p$ de $0{,}03$ es
  $P(\text{datos}\mid H_0)$. Leerlo como «hay 3\,\% de probabilidad de que mi
  hipótesis sea falsa» es $P(H_0\mid\text{datos})$: exactamente la misma
  inversión. Si entendiste este ejercicio, entendiste ése también.}
\enclase{El punto (b) es el que importa profesionalmente. Pide que lo lean en
voz alta: la mayoría escribe algo que sigue sonando a diagnóstico.}

\trampa{leer la sensibilidad como si respondiera la pregunta del paciente}%
  {los dos números son porcentajes altos sobre el mismo test, y la ficha
   técnica publica la sensibilidad pero no el valor predictivo}%
  {la sensibilidad divide por los que tienen la condición y el valor predictivo
   por los que dieron positivo. Para el paciente sólo sirve el segundo, y
   depende de la prevalencia de tu población}

\cierre{%
  El instrumento acierta el 88\,\% por los dos lados y aun así casi la mitad de
  sus alarmas son falsas. No es una contradicción: es lo que pasa cuando se
  busca algo poco frecuente, y se calcula con una división.
  \par
  Hasta aquí todas las preguntas fueron sobre eventos sueltos: dar positivo,
  tener el trastorno, presentar las dos sintomatologías. Lo que viene es
  describir el comportamiento completo de una variable de una sola vez.}
